\documentclass[11pt]{article}

\usepackage[a4paper,margin=1in]{geometry}
\usepackage{graphicx}
\usepackage{booktabs}
\usepackage{longtable}
\usepackage{hyperref}
\usepackage{enumitem}
\usepackage{amsmath}

\title{Does Orchestration Earn Its Cost?\\
A Comparative Study of Single-Agent vs. Multi-Agent LLM Systems\\
on Software Engineering Tasks}
\author{CS 540 Final Project Report\\
[Team Member Names]}
\date{May 2026}

\begin{document}

\maketitle
\tableofcontents
\newpage

\section{Introduction}

\subsection{Original Goals of the Project}

The original goal of this project was to compare single-agent and multi-agent large language model (LLM) systems on two canonical software engineering tasks: requirements engineering (RE) classification and code generation.

For requirements engineering, the system would receive a natural-language requirement sentence and classify it as:
\begin{itemize}
    \item Functional Requirement (FR)
    \item Non-Functional Requirement (NFR)
    \item Not a Requirement (NONE)
\end{itemize}

with an NFR subtype if applicable.

For code generation, the system would receive a programming problem and generate correct Python code, evaluated through automated test execution.

The central research question was whether multi-agent orchestration---decomposing a task across a Planner, Extractor, Critic, and Coder---improves output quality over a simpler single-shot approach, and whether any quality gain justifies the additional computational cost.

\subsection{Modified Goals}

The RE classification task was redesigned partway through the project into requirements elicitation. Rather than classifying pre-existing sentences, the system now receives a natural-language use-case description and generates a complete list of requirements from scratch.

This modification was made for two reasons:

\begin{enumerate}
    \item Classification of individual sentences is a low-complexity task where modern LLMs achieve near-ceiling performance in a single call, making it a poor testbed for comparing single-agent vs. multi-agent systems.
    
    \item Elicitation better reflects the actual challenge in requirements engineering practice: generating requirements from an informal description, not labeling pre-written ones.
\end{enumerate}

To support this redesign, use-case descriptions were synthetically generated from ground-truth requirement lists using the Claude API, giving a controlled input/output pair for evaluation.

The datasets changed from NICE (classification) and SecReq to NICE (elicitation) and PURE (elicitation).

A second modification was the addition of BigCodeBench-Hard as a third code generation benchmark. HumanEval+ and MBPP+ both produced 100\% pass@1 for single-agent and V1, making them too easy to reveal differences between systems.

BigCodeBench-Hard was added to stress-test the repair loop on real-world library-dependent tasks.

\subsection{Assumptions}

\begin{itemize}
    \item Platform: macOS (Darwin 25.0.0, Apple Silicon)
    \item Programming language: Python 3.14
    \item LLM: Anthropic \texttt{claude-sonnet-4-6}, temperature = 0
    \item Orchestration framework: LangGraph (LangChain)
    \item Evaluation:
    \begin{itemize}
        \item Code correctness assessed via subprocess execution
        \item RE quality assessed via semantic similarity
    \end{itemize}
    \item Cost model: Tokens estimated as:
\[
\frac{\text{input chars + output chars}}{4}
\]
    \item Reproducibility:
    \begin{itemize}
        \item temperature = 0
        \item RANDOM\_SEED = 42
    \end{itemize}
\end{itemize}

\subsection{Definitions}

\begin{description}
    \item[pass@1] Fraction of problems solved correctly on the first attempt.
    
    \item[Semantic F1] Harmonic mean of semantic precision and semantic recall using cosine similarity threshold 0.60.
    
    \item[Single-agent] One LLM call per task with limited retries on compile/parse failure.
    
    \item[Multi-agent V1] Planner $\rightarrow$ Extractor $\rightarrow$ Critic $\rightarrow$ Coder $\rightarrow$ Test Runner pipeline with repair loop.
    
    \item[Multi-agent V2] V1 extended with either:
    \begin{itemize}
        \item Test Critic node (CodeGen)
        \item SME advisory node (RE)
    \end{itemize}
    
    \item[Repair Loop] Iterative repair cycle triggered after failed tests.
    
    \item[Budget Guard] Hard constraints on:
    \begin{itemize}
        \item Maximum LLM calls
        \item Maximum tokens
    \end{itemize}
    
    \item[Context Trace] Per-invocation log tracking available context during each coder call.
\end{description}

\subsection{Summary of Approach}

We built three systems on the same underlying model and evaluated them under identical conditions.

For code generation:
\begin{itemize}
    \item HumanEval+
    \item MBPP+
    \item BigCodeBench-Hard
\end{itemize}

For requirements elicitation:
\begin{itemize}
    \item NICE
    \item PURE
\end{itemize}

After observing that all systems achieved 100\% pass@1 on HumanEval+ and MBPP+, we introduced BigCodeBench-Hard to expose regime differences.

During analysis, we identified two context-loss bugs in the V1 repair loop:
\begin{enumerate}
    \item Planner output dropped during repair
    \item Unittest-style tests were never executed
\end{enumerate}

Both bugs were fixed and the experiments were rerun.

\subsection{Project Deliverables}

\begin{enumerate}
    \item Three implemented systems:
    \begin{itemize}
        \item Single-agent
        \item Multi-agent V1
        \item Multi-agent V2
    \end{itemize}
    
    \item Five benchmark evaluations:
    \begin{itemize}
        \item HumanEval+
        \item MBPP+
        \item BigCodeBench-Hard
        \item NICE
        \item PURE
    \end{itemize}
    
    \item Context trace mechanism
    
    \item Bug analysis and fixes for context-loss issues
    
    \item Resume-safe evaluation scripts
    
    \item Results reports and benchmark summaries
\end{enumerate}

\subsection{Contribution of Individual Team Members}

[To be filled in]

\section{Background Research}

\subsection{Primary Research}

The following tools and frameworks were evaluated:

\begin{itemize}
    \item LangGraph
    \item Anthropic Claude API
    \item EvalPlus
    \item BigCodeBench-Hard
    \item sentence-transformers
\end{itemize}

\subsection{Secondary Research}

The following papers and resources informed the project:

\begin{enumerate}
    \item EvalPlus
    \item BigCodeBench
    \item NICE
    \item PURE
    \item Sentence-BERT
    \item LangGraph Documentation
\end{enumerate}

\subsection{What Was Most Helpful}

The BigCodeBench paper was the most influential because it revealed that easier benchmarks were insufficient for evaluating orchestration benefits.

LangGraph TypedDict documentation was critical for identifying silent state-loss issues inside the repair loop.

\section{Approach to Goals}

\subsection{Strategy}

The strategy was controlled ablation:
\begin{itemize}
    \item Hold model, prompts, and evaluation fixed
    \item Vary only orchestration architecture
\end{itemize}

Three systems were implemented:
\begin{enumerate}
    \item Single-agent baseline
    \item Multi-agent V1
    \item Multi-agent V2
\end{enumerate}

\subsection{Motivation for Strategy}

The controlled-ablation design isolates orchestration effects from:
\begin{itemize}
    \item Model quality
    \item Prompt engineering
    \item Framework differences
\end{itemize}

The study also compares two fundamentally different tasks:
\begin{itemize}
    \item Code generation (binary correctness)
    \item Requirements elicitation (semantic correctness)
\end{itemize}

\subsection{Division of Work Among Team Members}

[To be filled in]

\section{Implementation}

\subsection{Software Design}

The project is implemented as a Python package with the following modules:

\begin{itemize}
    \item \texttt{src/llm/client.py}
    \item \texttt{src/schemas/}
    \item \texttt{src/systems/single\_agent/}
    \item \texttt{src/systems/multi\_agent/nodes/}
    \item \texttt{src/evaluation/}
    \item \texttt{scripts/}
\end{itemize}

\subsection{External Components}

\begin{table}[h]
\centering
\begin{tabular}{lll}
\toprule
Component & Version & Purpose \\
\midrule
langgraph & latest & Multi-agent orchestration \\
anthropic & latest & Claude API access \\
evalplus & latest & Benchmark evaluation \\
sentence-transformers & latest & Semantic similarity \\
pydantic & v2 & Structured outputs \\
\bottomrule
\end{tabular}
\caption{External libraries and frameworks}
\end{table}

\subsection{Design of Empirical Studies}

\subsubsection*{Code Generation Evaluation}

Generated code was executed in isolated subprocesses with benchmark test suites.

\subsubsection*{Requirements Elicitation Evaluation}

Requirements were evaluated using sentence embeddings and cosine similarity matching.

\subsubsection*{Cost Tracking}

Token usage and LLM calls were tracked for all systems.

\section{Results}

\subsection{Empirical Results}

\subsubsection{Code Generation}

\begin{table}[h]
\centering
\begin{tabular}{lllll}
\toprule
System & Dataset & N & pass@1 & Avg Tokens \\
\midrule
Single-Agent & HumanEval+ & 164 & 100\% & 259 \\
Single-Agent & MBPP+ & 378 & 100\% & 138 \\
Single-Agent & BigCodeBench-Hard & 148 & 62.8\% & -- \\
Multi-Agent V1 & HumanEval+ & 164 & 100\% & 2037 \\
Multi-Agent V1 & MBPP+ & 378 & 100\% & 1531 \\
Multi-Agent V1 & BigCodeBench-Hard & 148 & 37.8\% & 5718 \\
\bottomrule
\end{tabular}
\caption{Code generation results}
\end{table}

\subsubsection{Requirements Elicitation}

\begin{table}[h]
\centering
\begin{tabular}{llll}
\toprule
Dataset & System & F1 & Coverage \\
\midrule
NICE & Single-Agent & 0.420 & 0.497 \\
NICE & Multi-Agent V1 & 0.348 & 0.424 \\
PURE & Single-Agent & 0.338 & 0.359 \\
PURE & Multi-Agent V1 & 0.325 & 0.329 \\
\bottomrule
\end{tabular}
\caption{Requirements elicitation results}
\end{table}

\subsubsection{Cross-Task Cost Analysis}

\begin{table}[h]
\centering
\begin{tabular}{llll}
\toprule
Task & Dataset & Token Overhead & Quality Change \\
\midrule
CodeGen & HumanEval+ & 7.9x & 0 \\
CodeGen & MBPP+ & 11.1x & 0 \\
RE & NICE & 1.9x & -17\% \\
RE & PURE & 1.7x & -4\% \\
\bottomrule
\end{tabular}
\caption{Cost-quality tradeoffs}
\end{table}

\subsection{Software Deliverables}

\begin{itemize}
    \item \texttt{scripts/run\_mbpp.py}
    \item \texttt{scripts/run\_bigcodebench.py}
    \item \texttt{scripts/run\_re\_elicitation.py}
    \item \texttt{scripts/evaluate\_re\_elicitation.py}
    \item \texttt{scripts/prepare\_re\_elicitation.py}
\end{itemize}

\section{Waiting Room}

The following items were planned but not completed:

\begin{enumerate}
    \item Fixed V1 full rerun on BigCodeBench-Hard
    \item Extended qualitative failure analysis
    \item Additional orchestration variants
\end{enumerate}

\bibliographystyle{plain}
\bibliography{references}

\end{document}
